%!TeX program=lualatex

\documentclass[12pt]{article}
\input{structure.tex}
\renewcommand{\bottomtitlespace}{0cm}

\begin{document}

\renewcommand{\tl}{$1$ сек.}
\renewcommand{\ml}{$1024$ MB}
\problem{Задача АB22. ФОРТУНА}
\addauthor{Автор: Емил Инджев}{2026}{03}{15}
\renewcommand{\header}{
	XLII НАЦИОНАЛНА ОЛИМПИАДА ПО ИНФОРМАТИКА\\
	Национален кръг\\
	Варна, 13 – 16 март 2026 г.\\
	Група AB, 9 – 12 клас, Ден 2\\
	\vspace{0.1cm}
}

Великият Цезар е изправен пред $M$-лема ($M$-странна дилема). След месеци неуспешни разсъждения, той решил да търси божествено вдъхновение в Храма на Фортуна (богинята на късмета, известна и като Автоматия). По-точно, той иска да получи произволно число между $0$ и $M-1$. Разбира се, числото трябва да е наистина произволно с точно равни вероятности за всички $M$ стойности. Това е невъзможно за простосмъртни и може да се случи само с помощта на богинята. За целта, в храма ще се извършва дневен ритуал, при който богинята осенява свой поклонник с едно наистина произволно число между $0$ и $N-1$ (с точно равни вероятности за всички $N$ стойности). За жалост на Цезар $N \neq M$. Затова той възложил на своите съветници задачата да измислят метод за произвеждане на търсения вид произволно число, използвайки редица числа от ритуала.

Това вече е доста трудна задача, но храмът си има допълнителни правила. Вътре се държи свитък с едно неотрицателно число $X$, записано на него (в началото това число е $0$ и е винаги по-малко от $500$). Всеки ден в храма нов поклонник извършва ритуала. Той вижда записаното число $X$ и получава произволното число $R$ от богинята. На база само тези две числа, той трябва да реши или да приключи целия процес и да обяви резултатното, търсено от Цезар, произволно число $T$, или, изтривайки старото, да запише ново неотрицателно число $Y$ в свитъка (отново $Y$ трябва да е по-малко от $500$).

Освен всичко това, Цезар иска средният брой дни за генериране на търсеното произволно число да е възможно най-малък. Забележете, че генерирането на произволното число може да отнеме произволно голям брой дни, стига това да се случва достатъчно рядко, че средният брой дни да е малък (изисква се той да е най-много $100$).

Помогнете на Цезар, като напишете програма \textbf{\texttt{fortuna}}, която имплементира такъв метод за генериране на произволно число от $0$ до $M-1$ използвайки произволни числа от $0$ до $N-1$.

\subsection{Детайли по имплементацията}

Трябва да имплементирате две функции:

\begin{lstlisting}
 void setup(int N, int M)
\end{lstlisting}

Тази функция ще се извика веднъж в началото с дадените $N$ и $M$.

\begin{lstlisting}
 std::pair<bool, int> proc(int X, int R)
\end{lstlisting}

Тази функция ще се извика по веднъж за всички възможни комбинации от $X$ и $R$, като възможни $X$ са само такива, които Вашето решение някога би записало в свитъка (и $0$). За всяка комбинация, функцията определя какво да се случи. За приключване и обявяване на $T$, функцията трябва да върне \verb|pair|-а \verb|{true, T}|. Зa продължаване и записване на $Y$ в свитъка, тя трябва да върне  \verb|pair|-а \verb|{false, Y}|.

\vspace{-0.1em}
След всички нужни извиквания на функциите Ви, програмата на журито ще провери, че Вашият алгоритъм произвежда всички $M$ възможни стойности за $T$ с точно равни вероятности и че никога не зацикля вечно (това се случва използвайки специални методи, които не са Ви предоставени). Освен това, програмата на журито ще симулира генериране на произволно число $10^6$ пъти и ще сметне средния брой дни $C$. Целта е да минимизирате $C$, като за $C > 100$ не се получават точки. 

\subsection{Ограничения}

\vspace{0.1em}

\begin{itemize}
\item $2 \leq N, M \leq 30$
\item $0 \leq X, Y < 500$
\item $0 \leq R < N$
\item $0 \leq T < M$
\end{itemize}

\subsection{Оценяване}

Всеки тест се оценява независимо от другите. Нека $C^*$ да е средният брой дни на оптималното възможно за теста решение. Тогава частта от точките, която ще получите за теста (ако решението Ви е валидно за него) ще е $S$:

$Q = \min(\frac{C^* + 0.005}{C}, 1)$

$S = \frac{1 + Q^{10} - {(1 - Q)}^{0.15}}{2}$

\subsection{Тестове}

\begin{itemize}
\item В първите $25\%$ от тестовете: $N > M$.
\item В следващите $25\%$ от тестовете: $N = 2$.
\item В последните $50\%$ от тестовете: няма допълнителни ограничения.
\end{itemize}

\subsection{Примерна комуникация}
\vspace{0.5 em}
\begin{table}[ht]
\begin{tblr}{width=0.85\textwidth,colspec={|Q[c]|Q[c]|}}
\hline
\textbf{Програма на журито} & \textbf{Вашата програма} \\
\hline
\texttt{setup(2, 3)} & \\
\hline
\texttt{proc(0, 0)} & \texttt{return \{false, 1\}} \\
\hline
\texttt{proc(0, 1)} & \texttt{return \{true, \ 0\}} \\
\hline
\texttt{proc(1, 0)} & \texttt{return \{false, 2\}} \\
\hline
\texttt{proc(1, 1)} & \texttt{return \{true, \ 1\}} \\
\hline
\texttt{proc(2, 0)} & \texttt{return \{false, 0\}} \\
\hline
\texttt{proc(2, 1)} & \texttt{return \{true, \ 2\}} \\
\hline
\end{tblr}
\end{table}
\FloatBarrier

Тук се иска да се реши задачата за $N=2$ и $M=3$. Един опит за решаване е описан в дадената примерна комуникация. Това решение прави ритуала до 3 пъти, докато не се падне единица и завършва с $T$ равно на номера на деня, в който първо се е паднала единица. Ако нито един от трите дни не се падне единица, то повтаря това отново, докато не успее. Това решение средно отнема $2$ дни, но за жалост не е вярно, защото не избира стойностите $0$, $1$ и $2$ за $T$ с равни верояности: $P(0) \approx 0.57$, $P(1) \approx 0.29$ и $P(2) \approx 0.14$.

\subsection{Локален грейдър}

Локалният грейдър не прави описаните проверки, а само симулира генериране на $10^6$ произволни числа и отпечатва $C$ и броят пъти, при които е била обявена всяка от $M$-те стойности за $T$. Локалният грейдър чете само $N$ и $M$ от стандартния вход (използва фиксиран сийд за рандомизация).

\end{document}
